\section{Introduction}
\label{nips:sec:intro}
\iffalse
DM \cite{COIFMAN20065,coifman2008diffusion}(coifman-lafon and the others) revolutionized/marked a turning point int non-linear dimension reduction, both wrt its theoretical understanding and its practical applicability. DM is now widely used in scientific data analysis \mmp{cite as many papers as we can} and is arguably the most studied dimension reduction method \mmp{cite as many as we can}. It also engendered a novel work in dimension reduction \cite{} and resurfaced/ XXXX traditional work \cite{}.

We now know that \mmp{again cite as many as we can, with facts}

This paper builds on the above body of knowledge to consider the following high level questions/distilling from it a ``profile'' for the DM embedding and its properties. 


ORIGINAL NEURIPS INTRO
\fi
Low dimensional embeddings of high-dimensional data are often used for visualization \cite{tsne}, i.e. as the {\em last step} of a data analysis. We ask whether non-linear dimension reduction, performed in a principled way, can also be useful as a {\em first step} of the data analysis, just like the classical Principal Component Analysis. Therefore, this paper examines to what extent the Diffusion Maps (DM) embedding\cite{COIFMAN20065,belkin2003laplacian}, a dimensionality reduction technique that captures the geometric structure of data by constructing a diffusion process among data points, obtains ``well conditioned'' non-linear embeddings, i.e. embeddings whose geometric distortions can be bounded.

Of particular interest is the preservation of the tangent space estimation under the DM embedding. For a smooth function on a manifold $M$, its gradient (as well as higher order derivatives) live in the tangent space to $M$. Controlling the error of the tangent space estimation in finite sample settings is a necessary condition to ensure that the gradients of functions of interest are transported unbiasedly to the lower dimensional data representation.

The DM is well studied as a differential geometric object; \cite{berard1994embedding} showed that in the limit of infinite dimension it is an isometric embedding, while \cite{portegies2016embeddings} showed that almost isometry can be achieved with a finite number $m$ of eigenfunctions, where this $m$ depends on manifold geometric properties. \cite{bates2014embedding} showed that if isometry is not required, then the sufficient embedding dimension by Laplacian eigenfunctions depends on dimension, injectivity radius, Ricci curvature and volume, thus that it can be arbitrarily larger than the Whitney embedding dimension of $2d$.

The DM embedding is widely used for non-linear dimension reduction as the Diffusion Maps algorithm \cite{COIFMAN20065,lafon2006diffusion,coifman2008diffusion,arvizu2016dimensionality}, which embeds a sample into $m$ dimensions by the eigenvectors of $L_n$, a $n\times n$ matrix estimator of the Laplace-Beltrami operator $\Delta$.
For instance, DM was used to analyze high-dimensional single-cell RNA sequencing data, revealing cell differentiation trajectories and underlying biological patterns\cite{haghverdi2015diffusion}, and it can order cells along their differentiation paths, enabling accurate reconstruction of branching developmental processes\cite{haghverdi2016diffusion}. 
In chemistry, DM can extract the dynamical modes of high-dimensional simulation trajectories, furnishing a kinetically low-dimensional framework\cite{ferguson2011nonlinear}, it can identify the collective coordinates of rare events in molecular transitions\cite{boninsegna2015investigating}.
In Astronomy, DM is applied to estimate galaxy redshifts from photometric data, demonstrating comparable accuracy to existing methods\cite{freeman2009photometric}.

Stimulated in part by the practical applications of DM, the consistency of  finite sample estimators of $\Delta$, and of their eigenfunctions has been intensely studied. Many early works discussed the pointwise consistency for empirical operator with smooth functions on $M$ and they derived convergence (with different rates) \cite{von2008consistency,belkin2008towards,gine2006empirical,10.5555/3104322.3104459,hein2007graph}, which provide illuminating results for subsequent studies. Not only $\Delta$ itself, the consistency of its embedding are also widely studied, they discussed how the eigenvectors and eigenvalues of the empirical Laplacian with specific graph (K-NN, $\epsilon$-graph) converge to the eigenfunctions and eigenvalues of $\Delta$ in the different norm ($L^2$, $L^\infty$)\cite{cheng2022eigen,garcia2020error,wang2015spectral,belkin2006convergence}. There are some other special convergence types studied. For instance, convergence in Lipschitz norm\cite{calder2022lipschitz}, embedding using heat kernel\cite{jones2008manifold,DUNSON2021282}

This paper completes the picture by considering geometric properties of a manifold, such as smoothness, injectivity radius, reach, volume, diameter, and examining to what extent the DM preserves these properties. In other words, if the original manifold $\MM$ is ``well behaved'', what can be said about its DM embedding $\varphi(\MM)$? Surprisingly, these questions have received less attention in the literature.

Specifically, under a set of mostly standard geometric assumptions on a family of compact manifolds (\cref{sec:problem}), in \cref{sec:phim-L} we study the smoothness of $\varphi$ by examining its Sobolev norms (\cref{sec:smooth}) and use it to show the existence of local parameterizations, then we derived consistent properties that hold uniformly for all manifolds in the family after DM embedding, like sampling densities and reach (\cref{phiM}). In \cref{sec:convergence}, we identify the correspondence between the embedding via eigenfunctions and the embedding via eigenvectors of the graph matrix and address the convergence between them, then estimate the error between the eigenfunctions and eigenvectors to quantify the noise in the positions of the embedded points. \cref{sec:experiments} presents the experimental results, and \cref{sec:conclusion} concludes the paper.

